%% abtex2-modelo-trabalho-academico.tex, v-1.9.2 laurocesar
%% Copyright 2012-2014 by abnTeX2 group at http://abntex2.googlecode.com/ 
%%
%% This work may be distributed and/or modified under the
%% conditions of the LaTeX Project Public License, either version 1.3
%% of this license or (at your option) any later version.
%% The latest version of this license is in
%%   http://www.latex-project.org/lppl.txt
%% and version 1.3 or later is part of all distributions of LaTeX
%% version 2005/12/01 or later.
%%
%% This work has the LPPL maintenance status `maintained'.
%% 
%% The Current Maintainer of this work is the abnTeX2 team, led
%% by Lauro César Araujo. Further information are available on 
%% http://abntex2.googlecode.com/
%%
%% This work consists of the files abntex2-modelo-trabalho-academico.tex,
%% abntex2-modelo-include-comandos and abntex2-modelo-references.bib
%%

% ------------------------------------------------------------------------
% ------------------------------------------------------------------------
% abnTeX2: Modelo de Trabalho Academico (tese de doutorado, dissertacao de
% mestrado e trabalhos monograficos em geral) em conformidade com 
% ABNT NBR 14724:2011: Informacao e documentacao - Trabalhos academicos -
% Apresentacao
% MODELO MODIFICADO POR JHON WILKER DA S. SOUSA PARA ATENDER AS ESPECIFICIDADES DO CURSO DE CIÊNCIA DA COMPUTAÇÃO DA UNIVERSIDADE ESTADUAL DO PIAUI
% ------------------------------------------------------------------------
% 

\documentclass[
	% -- opções da classe memoir --
	12pt,				% tamanho da fonte
	openright,			% capítulos começam em pág. ímpar (insere página vazia caso preciso)
	oneside,			% para impressão em verso e anverso. Oposto a oneside
	a4paper,			% tamanho do papel. 
	% -- opções da classe abntex2 --
	chapter=TITLE,		% títulos de capítulos convertidos em letras maiúsculas
	%section=TITLE,		% títulos de seções convertidos em letras maiúsculas
	%subsection=TITLE,	% títulos de subseções convertidos em letras maiúsculas
	%subsubsection=TITLE,% títulos de subsubseções convertidos em letras maiúsculas
	% -- opções do pacote babel --
	english,			% idioma adicional para hifenização
	%french,				% idioma adicional para hifenização
	%spanish,			% idioma adicional para hifenização
	brazil,				% o último idioma é o principal do documento
	]{abntex2}
% ---
% Pacotes básicos 
% ---
%\usepackage{lmodern}			% Usa a fonte Latin Modern	
\usepackage{mathptmx}
\renewcommand{\ABNTEXchapterfont}{\normalfont}
%\usepackage{bookman}		
\usepackage[T1]{fontenc}		% Selecao de codigos de fonte.
\usepackage[utf8]{inputenc}		% Codificacao do documento (conversão automática dos acentos)
\usepackage{lastpage}			% Usado pela Ficha catalográfica
\usepackage{indentfirst}		% Indenta o primeiro parágrafo de cada seção.
\usepackage{color}				% Controle das cores
\usepackage{graphicx}			% Inclusão de gráficos
\usepackage{microtype} 			% para melhorias de justificação


\usepackage{epigraph}

%Depois no documento escrevemos:

%\epigraph{texto}{referência}


%\usepackage{geometry}
%\geometry{a4paper, left=3cm, right=2cm, bottom=2cm, top=3cm}
% ---
\graphicspath{ {./figs/} }
% ---
% Pacotes adicionais, usados apenas no âmbito do Modelo Canônico do abnteX2
% ---
\usepackage{lipsum}				% para geração de dummy text
% ---
% Pacotes de citações
% ---
\usepackage[brazilian,hyperpageref]{backref}	 % Paginas com as citações na bibl
\usepackage[alf]{abntex2cite}	% Citações padrão ABNT

% --- 
% CONFIGURAÇÕES DE PACOTES
% --- 

% ---
% Configurações do pacote backref
% Usado sem a opção hyperpageref de backref
\renewcommand{\backrefpagesname}{Citado na(s) página(s):~}
% Texto padrão antes do número das páginas
\renewcommand{\backref}{}
% Define os textos da citação
\renewcommand*{\backrefalt}[4]{
	\ifcase #1 %
		Nenhuma citação no texto.%
	\or
		Citado na página #2.%
	\else
		Citado #1 vezes nas páginas #2.%
	\fi}%
% ---

% ---
% Informações de dados para CAPA e FOLHA DE ROSTO
% ---
\titulo{TITULO DO SEU TRABALHO}
\autor{Nome do Aluno Completo}
\local{Teresina}
\data{2023}
\orientador{Nome do Orientador Completo}
\instituicao{Universidade Estadual do Piauí}
\tipotrabalho{Monografia (Graduação)}
% O preambulo deve conter o tipo do trabalho, o objetivo, o nome da instituição e a área de concentração 
\preambulo{Monografia submetida à Universidade Estadual do Piauí como parte dos requisitos para a obtenção do grau de Bacharel em Ciência da Computação.}




% Configurações de aparência do PDF final
%\definecolor{blue}{RGB}{41,5,195}
% informações do PDF
\makeatletter
\hypersetup{
     	%pagebackref=true,
		pdftitle={\@title}, 
		pdfauthor={\@author},
    	pdfsubject={\imprimirpreambulo},
	    pdfcreator={LaTeX with abnTeX2},
		pdfkeywords={abnt}{latex}{abntex}{abntex2}{trabalho acadêmico}, 
		colorlinks=true,       		% false: boxed links; true:              colored links
    	linkcolor=black,          	% color of internal links
    	citecolor=black,        		% color of links to bibliography
    	filecolor=magenta,      		% color of file links
		urlcolor=black,
		bookmarksdepth=4}

\makeatother
% --- 
% Espaçamentos entre linhas e parágrafos 
% --- 
% O tamanho do parágrafo é dado por:
\setlength{\parindent}{1.3cm}
% Controle do espaçamento entre um parágrafo e outro:
\setlength{\parskip}{0.2cm}  % tente também \onelineskip
% ---
% compila o indice
% ---
\makeindex
% ---

% --------------------------------------------------------------------------------------------------------
% % Início do documento
% --------------------------------------------------------------------------------------------------------
\begin{document}

% Retira espaço extra obsoleto entre as frases.
\frenchspacing 

% Capa
%---------------------------------------------------------------------------------------------------------

    \begin{center}
    {\ABNTEXchapterfont\large UNIVERSIDADE ESTADUAL DO PIAUÍ}\\
    {\ABNTEXchapterfont\large CIÊNCIA DA COMPUTAÇÃO}
    
    \vspace*{\fill}\vspace*{\fill}
    {\ABNTEXchapterfont\large \imprimirautor}  
    \vspace*{\fill}
   
    \vspace*{\fill}\vspace*{\fill}
    \begin{center}
    \ABNTEXchapterfont \bfseries \Large \imprimirtitulo
    \end{center}
    \vspace*{\fill}\vspace*{\fill}
   
  
   \end{center}  

      
    \begin{center}
    \vspace*{0.5cm}
    {\ABNTEXchapterfont\large TERESINA}
    \par
    {\ABNTEXchapterfont\large \imprimirdata}
    \vspace*{1cm}
    \end{center}
%--------------------------------------------------x------------------------------------------------------




% Folha de Rosto
% --------------------------------------------------------------------------------------------------------
\begin{folhaderosto}

  \begin{center}
    {\ABNTEXchapterfont\large \imprimirautor}

    \vspace*{\fill}\vspace*{\fill}
    \begin{center}
      \ABNTEXchapterfont\bfseries\Large \imprimirtitulo
    \end{center}
    \vspace*{\fill}
    
    \hspace{.45\textwidth}
    \begin{minipage}{.5\textwidth}
        \imprimirpreambulo
    \end{minipage}%
    \vspace{1.5cm}
    Orientador: \imprimirorientador \\
    
    \vspace*{\fill}
     
   \end{center}  

     \begin{center}
    \vspace*{0.5cm}
    {\ABNTEXchapterfont\large TERESINA}
    \par
    {\ABNTEXchapterfont\large \imprimirdata}
    \vspace*{1cm}
    \end{center}
  
\end{folhaderosto}
%--------------------------------------------------x------------------------------------------------------

% Folha ficha Catalográfica
% --------------------------------------------------------------------------------------------------------

%Substitua pela ficha catalografica gerada pela biblioteca da UESPI
\begin{fichacatalografica}
	\vspace*{\fill}					% Posição vertical
	\hrule							% Linha horizontal
	\begin{center}					% Minipage Centralizado
	\begin{minipage}[c]{12.5cm}		% Largura
	
	\imprimirautor
	
	\hspace{0.5cm} 
	{\imprimirtitulo / \imprimirautor. --	\imprimirlocal~~-- PI, Brasil, \imprimirdata.}
	
	{\hspace{0.5cm}\imprimirorientadorRotulo~\imprimirorientador}
	
	\hspace{0.5cm}\parbox[t]{\textwidth}{\imprimirtipotrabalho~--~\imprimirinstituicao, \imprimirdata.}

    \hspace{0.5cm}
		1. Embarcado.
		2. GSM.
        3. GPS.
        3. RF 433 Mhz.
		I. Nome do Orientador.
		II. Universidade Estadual do Piauí.
		 			
%\hspace{8.75cm} CDU 02:141:005.7\\
	
	\end{minipage}
	\end{center}
	\hrule
\end{fichacatalografica}
%--------------------------------------------------x------------------------------------------------------

% Folha de Aprovação
% --------------------------------------------------------------------------------------------------------

\begin{folhadeaprovacao}

  \begin{center}
    {\ABNTEXchapterfont\bfseries\Large \imprimirtitulo}\\
    \vspace*{\fill}
    {\ABNTEXchapterfont\large \imprimirautor}
    
    \vspace*{\fill}\vspace*{\fill}
    \hspace{.45\textwidth}
    \begin{minipage}{.5\textwidth}
        \imprimirpreambulo
    \end{minipage}%
    \vspace*{\fill}
   
   \ABNTEXchapterfont     
   {Trabalho aprovado. Teresina, XX de Junho de \imprimirdata}
   \end{center}
   

   \assinatura{\textbf{\imprimirorientador} \\ Orientador} 
   \assinatura{\textbf{Nome do Avaliador-1 Completo} \\ }
   \assinatura{\textbf{Nome do Avaliador-2 Completo} \\ }

  
\end{folhadeaprovacao}
%--------------------------------------------------x------------------------------------------------------


% Dedicatória
%---------------------------------------------------------------------------------------------------------
\include{1_pretextuais/dedicatoria}
%--------------------------------------------------x------------------------------------------------------

% Agradecimentos
% --------------------------------------------------------------------------------------------------------
\include{1_pretextuais/agradecimentos}
%--------------------------------------------------x------------------------------------------------------

% Epígrafe
% --------------------------------------------------------------------------------------------------------
\include{1_pretextuais/epigrafe}
%--------------------------------------------------x------------------------------------------------------


% Resumo
%---------------------------------------------------------------------------------------------------------
\setlength{\absparsep}{18pt} % ajusta o espaçamento dos parágrafos do resumo
\begin{resumo}
\ABNTEXchapterfont
Consiste na apresentação dos pontos relevantes de um texto. O resumo deve dar uma visão rápida e clara do trabalho; constitui-se em uma sequencia de frases concisas e objetivas e não de uma simples enumeração de tópicos. Apresenta os objetivos do estudo, o problema, a metodologia, resultados alcançados e conclusão. Deve ser digitado em espaço simples e sem parágrafos, não ultrapassando a 500 palavras.


\textbf{Palavras-chaves}: Palavra-1, Palavra-2, Palavra-3.
\end{resumo}
%--------------------------------------------------x------------------------------------------------------


% Abstract
%---------------------------------------------------------------------------------------------------------
\begin{resumo}[Abstract]
\ABNTEXchapterfont
Consiste em traduzir o texto do resumo para o idioma ingles.

% \begin{otherlanguage*}{english}
%   This is the english abstract.
%   \vspace{\onelineskip}
%   \noindent 

\textbf{Keywords}: word1. Word-2. Word-3.
\end{resumo}
%--------------------------------------------------x------------------------------------------------------


% inserir lista de ilustrações
%---------------------------------------------------------------------------------------------------------
\pdfbookmark[0]{\listfigurename}{lof}
\listoffigures*
\cleardoublepage
%--------------------------------------------------x------------------------------------------------------


% inserir lista de tabelas
%---------------------------------------------------------------------------------------------------------
\pdfbookmark[0]{\listtablename}{lot}
\listoftables*
\cleardoublepage
%--------------------------------------------------x------------------------------------------------------


% inserir lista de abreviaturas e siglas
%---------------------------------------------------------------------------------------------------------
\begin{siglas}
\ABNTEXchapterfont\item[TCP] Protocolo de Controle de Transmissão
\item[UDP] Protocolo de Datagrama do Usuário 
\item[VoIP] Voz sobre IP 
\item[RTP] Protocolo Transporte em Tempo Real
\item[MPRTP] Protocolo de Transporte em Tempo Real de Caminho Multiplos \item[RTCP] Protocolo de Controle de Transporte em Tempo Real
\item[MPTCP] Protocolo de Controle de Transporte de Caminho Multiplo 
\end{siglas}
%--------------------------------------------------x------------------------------------------------------
\pdfbookmark[0]{\contentsname}{toc}
\tableofcontents*
\cleardoublepage
\textual
%--------------------------------------------------x------------------------------------------------------
\input{2_capitulos_tex/1_Introducao}
%--------------------------------------------------x------------------------------------------------------
\input{2_capitulos_tex/2_Referencial}
 % Desenvolvimento %
%---------------------------------------------------------------------------------------------------------
\input{2_capitulos_tex/3_Desenvolvimento}
%--------------------------------------------------x------------------------------------------------------

 % Conclusões e Trabalhos Futuros %
%---------------------------------------------------------------------------------------------------------
\input{2_capitulos_tex/4_Conclusao}
%--------------------------------------------------x------------------------------------------------------

%--------------------------------------------------x-----------------------------------------------------

%--------------------------------------------------x-----------------------------------------------------


%--------------------------------------------------x-----------------------------------------------------
% ----------------------------------------------------------
% ELEMENTOS PÓS-TEXTUAIS
% ----------------------------------------------------------
%{Referencias Bibliograficas}
%---------------------------------------------------------------------------------------------------------
\bibliography{biblio.bib}
\printindex

% ---------------------------------------------------------------------
    % GLOSSÁRIO
    % --------------------------------------------------------------------- 
    % Arquivo que contém as definições que vão aparecer no glossário
    %\input{tex/glossario}
    % Comando para incluir todas as definições do arquivo glossario.tex
    %\glsaddall
    % Impressão do glossário
    %\printglossaries

    % ----------------------------------------------------------
    % Apêndices
    % ----------------------------------------------------------
    
    % ---
    % Inicia os apêndices
    % ---
   % \apendices{
    %    \chapter{References Set}
    %    \label{apendice:setinicial}
    %    \input{trabRelacionados}

   % }
    % ---
    

    % ----------------------------------------------------------
    % Anexos
    % ----------------------------------------------------------
    
    % ---
    % Inicia os anexos
    % ---
    %\anexos{
    
    
    %\begin{anexosenv}
    
        
    %\end{anexosenv}
    %}
    % ---


\end{document}